\chapter{Sea-Wick Contractions and Band-Edge Corrections}
\label{app:interactions_bogoliubov}

\section{Scope and notation}

This appendix records the finite-band CAR bookkeeping behind the sea-Wick contraction and the diagonal correction used in Chapter~\ref{ch:luttinger_models}. The physical distinction between hopping and pairing, and the Bogoliubov transformation itself, are developed in the main chapters. Here we retain the combinatorial details that are useful when deriving the current interaction from the raw quartic expression.

\section{Syntactic Normal Ordering and the Quartic Contraction Identity}

Let $h \ge 1$ and let $L = 2h$ be the even lattice size. The momentum band is $B = \{-h+1, \dots, h\}$. The Dirac sea reference configuration is $S_0 = \{k \in B \mid k \le 0\}$, with sea indicator $s_k = 1$ if $k \le 0$ and $s_k = 0$ if $k > 0$.

\begin{definition}[Sea Quasiparticle Words and Normal Ordering]
For each momentum mode $k \in B$, we associate a sea quasiparticle type:
\begin{itemize}
    \item A creation letter $c_k^\dagger$ corresponds to a sea quasi-creator if $k > 0$, and to a sea quasi-annihilator if $k \le 0$.
    \item An annihilation letter $c_k$ corresponds to a sea quasi-annihilator if $k > 0$, and to a sea quasi-creator if $k \le 0$.
\end{itemize}
A word $w = (l_1, \dots, l_n)$ is an ordered list of labeled letters. The \emph{sea normal-ordered word} $:w:$ is the reordered product where all quasi-creators appear to the left of all quasi-annihilators, maintaining their relative stable order within each group, multiplied by the fermionic sign $(-1)^{\mathrm{inv}(w)}$, where $\mathrm{inv}(w)$ is the number of inversions required to sort quasi-creators past quasi-annihilators.
\end{definition}

Evaluating this syntactic definition on bilinear words yields the sea-normal-ordered bilinear operator:
\begin{equation}
:\! c_a^\dagger c_b \!: \;\coloneqq c_a^\dagger c_b - \delta_{a, b} s_a I.
\end{equation}

\begin{lemma}[Diagonal Bilinear Sea Ordering]
\label{lem:app_sea_normal_bilinear_diag}
For any mode $k \in B$ and species $\nu$:
\begin{equation}
:\! c_{\nu, k}^\dagger c_{\nu, k} \!: \;= \begin{cases} c_{\nu, k}^\dagger c_{\nu, k}, & \text{if } k > 0, \\ - c_{\nu, k} c_{\nu, k}^\dagger, & \text{if } k \le 0. \end{cases}
\end{equation}
\end{lemma}

\begin{proof}
By the canonical anticommutation relations, $c_{\nu, k}^\dagger c_{\nu, k} + c_{\nu, k} c_{\nu, k}^\dagger = I$. For $k \le 0$, $s_k = 1$, so $c_{\nu, k}^\dagger c_{\nu, k} - I = -c_{\nu, k} c_{\nu, k}^\dagger$. For $k > 0$, $s_k = 0$, so $c_{\nu, k}^\dagger c_{\nu, k} - 0 = c_{\nu, k}^\dagger c_{\nu, k}$.
\end{proof}

\begin{lemma}[Sea-Ordered Density Modes]
\label{lem:app_sea_normal_rho}
For any transfer $m \in \mathbb{Z}$, define the sea-ordered density mode by:
\begin{equation}
:\! \rho_{m, \nu} \!: \;\coloneqq \sum_{(p, k) \in D_m} :\! c_{\nu, p}^\dagger c_{\nu, k} \!:.
\end{equation}
Then:
\begin{enumerate}
    \item If $m \neq 0$, the transfer is non-zero, so $p = k + m \neq k$. Hence $\delta_{p, k} = 0$, and:
    \begin{equation}
    :\! \rho_{m, \nu} \!: \;= \rho_{m, \nu} = \sum_{(p, k) \in D_m} c_{\nu, p}^\dagger c_{\nu, k}.
    \end{equation}
    \item If $m = 0$, the transfer vanishes, so $p = k$. Then $:\! \rho_{0, \nu} \!:$ is identically the total charge operator:
    \begin{equation}
    :\! \rho_{0, \nu} \!: \;= \sum_{k \in B} :\! c_{\nu, k}^\dagger c_{\nu, k} \!: \;= \hat{N}_\nu.
    \end{equation}
\end{enumerate}
\end{lemma}

\begin{proof}
Part (1) follows because every pair $(p, k) \in D_m$ with $m \neq 0$ satisfies $p - k = m \neq 0$, so the diagonal subtraction $\delta_{p, k} s_k I$ vanishes identically.
Part (2) follows by summing Lemma~\ref{lem:app_sea_normal_bilinear_diag} over all $k \in B$. On any occupation ket $|S\rangle$, $\sum_{k \in B} c_{\nu, k}^\dagger c_{\nu, k} |S\rangle = |S_\nu| |S\rangle$, and $\sum_{k \in B} s_k = h$. Hence the sum evaluates to $(|S_\nu| - h) |S\rangle = N_\nu |S\rangle = \hat{N}_\nu |S\rangle$.
\end{proof}

\begin{theorem}[Exact Quartic Sea-Wick Contraction Identity]
\label{thm:app_quartic_sea_wick}
Let $m \in \mathbb{Z}$, and let $(a, b) \in D_m$ and $(c, d) \in D_{-m}$ be valid pairs in $B$, so that $a = b + m$ and $c = d - m$.
Then on the full finite Fock space $\mathcal{F}_\sigma(L)$:
\begin{equation}
\label{eq:app_quartic_wick}
:\! c_a^\dagger c_b c_c^\dagger c_d \!: \;= \;:\! c_a^\dagger c_b \!:\, :\! c_c^\dagger c_d \!: \;+ \; \delta_{d, a} \Big[ s_d t_b - (1 - s_b) t_d - s_d (1 - s_b) I \Big],
\end{equation}
where $t_k \coloneqq c_k^\dagger c_k - s_k I$.
Furthermore, on repeated single-mode indices ($m = 0$ and $a = b = c = d$):
\begin{equation}
:\! c_a^\dagger c_a c_a^\dagger c_a \!: \;= 0.
\end{equation}
\end{theorem}

\begin{proof}
The proof evaluates both sides on all $2^4 = 16$ possible combinations of sea indicators $s_a, s_b, s_c, s_d \in \{0, 1\}$.
In each case, one expresses the letters in terms of $q_k, q_k^\dagger$, applies the CAR anticommutation relations $q_i q_j^\dagger = \delta_{ij} I - q_j^\dagger q_i$, and compares the normally ordered word with the product of bilinears.
When $d = a$, moving the inner creation operator $c_c^\dagger = c_{a-m}^\dagger$ past the annihilation operator $c_b = c_{a-m}$ produces a non-vanishing contraction $\delta_{b, c} I = \delta_{a, d} I$. Tracking the signs and subtracting the product of bilinear vacuum expectations yields precisely the contraction bracket $s_d t_b - (1 - s_b) t_d - s_d(1 - s_b) I$.
For repeated indices $a = b = c = d$, $m = 0$ and $d = a$, so the right-hand side is $t_a^2 + (2s_a - 1)t_a$. Since $n_a^2 = n_a$, direct evaluation shows $t_a^2 + (2s_a - 1)t_a = 0$, matching the fermionic nilpotence $(c_a^\dagger)^2 = 0$.
\end{proof}

\begin{remark}[Formalization Note]
The formal statement and proof use the finite mathematical carriers and hypotheses specified in this chapter. The result is an exact algebraic or combinatorial claim within that scope; it does not presume an analytic limit or a stronger continuum identity.
\end{remark}

\section{Summing Transfers: The Diagonal One-Body Correction}

Next, we sum the quartic Wick identity over all transfers $m \in [-M, M]$ and all valid pairs $(a, b) \in D_m$ and $(c, d) \in D_{-m}$ for a single species $\nu$.

\begin{theorem}[Species Raw vs.\ Current Equivalence]
\label{thm:app_species_raw_current}
Let $M \ge 0$ and let $\nu \in \sigma$. The raw intra-branch sum satisfies:
\begin{equation}
\sum_{m=-M}^M \sum_{(a, b) \in D_m} \sum_{(c, d) \in D_{-m}} :\! c_{\nu, a}^\dagger c_{\nu, b} c_{\nu, c}^\dagger c_{\nu, d} \!: \;= 2 \sum_{m=1}^M \rho_{m, \nu} \rho_{-m, \nu} + \hat{N}_\nu^2 - Q_{M, \nu},
\end{equation}
where $Q_{M, \nu}$ is the diagonal one-body correction operator:
\begin{equation}
Q_{M, \nu} \coloneqq \sum_{k \le 0} \Big( 1 + 2\min(M, -k) \Big) (I - n_{\nu k}) + \sum_{k > 0} \Big( 1 + 2\min(M, k-1) \Big) n_{\nu k}.
\end{equation}
\end{theorem}

\begin{proof}
By Theorem~\ref{thm:app_quartic_sea_wick}, the sum splits into a product-of-bilinears sum and a contraction sum:
\[
\sum_{m=-M}^M \left( \sum_{(a, b) \in D_m} :\! c_{\nu, a}^\dagger c_{\nu, b} \!: \right) \left( \sum_{(c, d) \in D_{-m}} :\! c_{\nu, c}^\dagger c_{\nu, d} \!: \right) + \sum_{m=-M}^M \sum_{(a, b) \in D_m} \Big[ s_a t_b - (1-s_b)t_a - s_a(1-s_b)I \Big].
\]
The first sum is $\sum_{m=-M}^M :\! \rho_{m, \nu} \!:\, :\! \rho_{-m, \nu} \!:$. Isolating $m = 0$ gives $:\!\rho_{0, \nu}\!:^2 = \hat{N}_\nu^2$. For $m \neq 0$, $:\!\rho_{m, \nu}\!: = \rho_{m, \nu}$. Pairing positive and negative modes:
\[
\rho_{m, \nu} \rho_{-m, \nu} + \rho_{-m, \nu} \rho_{m, \nu} = 2 \rho_{m, \nu} \rho_{-m, \nu} + [\rho_{-m, \nu}, \rho_{m, \nu}].
\]
The finite-band edge commutator $[\rho_{-m, \nu}, \rho_{m, \nu}] = \sum_{(a, b) \in D_m} (n_{\nu b} - n_{\nu a})$ precisely combines with the contraction terms.
Summing the net indicators across all accessible transfers $n \in [1, M]$ for each fixed $k \in B$ yields the weight $1 + 2\min(M, -k)$ for holes ($k \le 0$) and $1 + 2\min(M, k-1)$ for particles ($k > 0$).
\end{proof}

\begin{lemma}[Properties of the Sea-Wick Correction $Q_{M, \nu}$]
\label{lem:app_q_properties}
\begin{enumerate}
    \item \textbf{Vacuum action:} On the zero-charge Dirac sea $|\Omega\rangle = |\vec{0}\rangle_0$:
    \begin{equation}
    Q_{M, \nu} |\Omega\rangle = 0.
    \end{equation}
    \item \textbf{Saturation of the diagonal correction:} If $M \ge h - 1$, the minima appearing in the weights of $Q_{M,\nu}$ have reached their band-dependent values. Then $\min(M, -k) = -k$ for $k \le 0$, and $\min(M, k-1) = k-1$ for $k > 0$. Consequently:
    \begin{equation}
    Q_{M, \nu} = 2\hat{P}_\nu - \hat{N}_\nu,
    \end{equation}
    where $\hat{P}_\nu = \sum_{k \in B} k (n_{\nu k} - s_k I)$ is the computational momentum.
\end{enumerate}
\end{lemma}

\begin{proof}
For (1), on $|\Omega\rangle$, $n_{\nu k} = 1$ for all $k \le 0$ (so $I - n_{\nu k} = 0$), and $n_{\nu k} = 0$ for all $k > 0$. Hence every summand vanishes.
For (2), substitution of the saturated minima into the definition of $Q_{M, \nu}$ yields $\sum_{k \le 0} (1 - 2k)(I - n_{\nu k}) + \sum_{k > 0} (2k - 1)n_{\nu k}$. Expressing this in terms of $t_k = n_{\nu k} - s_k I$ gives exactly $2 \sum_{k \in B} k t_k - \sum_{k \in B} t_k = 2\hat{P}_\nu - \hat{N}_\nu$.
\end{proof}
